% \iffalse meta-comment
%
% hit-ist.dtx 是索引的排序样式
%
% \fi
%
% \section{索引排序样式}
%
% \changes{v3.2a}{2026/08/09}{索引排序样式从 \file{hithesis.dtx} 单独成文件}
% 生成 \file{hithesis.ist}，是 \texttt{makeindex} 的配置，不是 \hologo{LaTeX} 代码。
% 跟 \file{hit-bst.dtx} 一样，都是喂给外部工具的，所以单独一个文件，不进 pages。
%
% 我工要求的索引格式，照规范附录 6 的官方示例（iota-hit 逐条量过）：
% 分组标题是单个大写字母（\texttt{headings\_flag 1}），居中、加粗、小四，
% 组间段前 12\,bp（\cs{addvspace}，首组贴着章标题的段后距被吸收，正与
% 示例“首组没有”一致）；组间不再另发 \cs{indexspace}（\texttt{group\_skip}
% 收成一个换行）；页码右对齐到栏缘（\texttt{delim} 一律 \cs{hfill}）。
% 条目字号跟正文小四，原先收尾的 \cs{wuhao} 去掉。
% \changes{v1.0.10}{2018/02/19}{修改了索引的间距，使其更符合规范中的示例}
%    \begin{macrocode}
%<*ist>
headings_flag 1
heading_prefix "\{\\par\\addvspace\{12bp\}\\centering\\normalsize\\textbf\{"
heading_suffix "\}\\par\}\\nopagebreak\n"
group_skip "\n"
delim_0 "\\hspace*{\\fill}"
delim_1 "\\hspace*{\\fill}"
delim_2 "\\hspace*{\\fill}"
%</ist>
%    \end{macrocode}
%
% \Finale
